% ============================================================
% theme_academique.tex — DA reprise de "Dossier scientifique"
%
% Palette de couleurs, barres de titre (section/sous-section/.../
% sous-paragraphe) et boîtes à bande colorée (\EncadreUn, \EncadreDeux)
% reprises quasi telles quelles de tex/commands.tex du Dossier
% scientifique fourni. Les boîtes pédagogiques du projet (coup de pouce,
% correction, solution détaillée, remarque, aller plus loin) sont
% redéfinies avec ce même langage visuel plutôt qu'avec l'ancien style
% "titre plein + fond teinté" (conservé dans theme_simple.tex).
%
% Nécessite (déjà chargés par packages.tex) : tikz, tcolorbox[most],
% et en plus ici : titlesec, varwidth.
% ============================================================
\usepackage{titlesec}
\usepackage{varwidth}

% ------------------------------------------------------------------
% Palette (reprise du Dossier scientifique)
% ------------------------------------------------------------------
\definecolor{colorStitre}{RGB}{140,0,0}       % section    -- rouge sombre
\definecolor{colorSStitre}{RGB}{190,140,140}  % sous-section -- rose poussiéreux
\definecolor{colorSSStitre}{RGB}{140,140,140} % sous-sous-section -- gris

\definecolor{Prune}{RGB}{99,0,60}
\definecolor{woodNoyer}{HTML}{3A2A1F}
\definecolor{woodChene}{HTML}{4A3628}
\definecolor{greenForet}{HTML}{0B3D2E}
\definecolor{greenSapin}{HTML}{0F4C3A}
\definecolor{grayPierre}{HTML}{4A4A4A}
\definecolor{grayBrume}{HTML}{7A7A7A}
\definecolor{headerblue}{HTML}{1565C0}
\definecolor{colorOchre}{HTML}{C98A2C}        % ajout : accent chaud pour le "coup de pouce"

% ============================================================
% COULEURS DES BLOCS PÉDAGOGIQUES
% ============================================================

% Prérequis : couleur terre / brun chaud
\definecolor{colorPrerequis}{HTML}{8B5E3C}

% Objectifs pédagogiques étudiants : bleu pédagogique
\definecolor{colorObjectifsPedagogiquesEtudiants}{HTML}{1565C0}

% Objectifs pédagogiques enseignants : prune / violet institutionnel
\definecolor{colorObjectifsPedagogiquesEnseignants}{HTML}{63365C}

% Point de cours : vert forêt
\definecolor{colorPointCours}{HTML}{0F4C3A}

% Remarques pédagogiques / écueils classiques : ocre
\definecolor{colorRemarquesPedagogiquesEcueilsClassiques}{HTML}{C98A2C}

% Remarques pédagogiques / conseils d'animation : bleu-gris
\definecolor{colorRemarquesPedagogiquesConseilsAnimation}{HTML}{57659E}

% Bilan : gris anthracite
\definecolor{colorBilan}{HTML}{4A4A4A}

% Palette numérotée (colorOne..colorSix) pour les annotations \ptFleche
% libres (voir templates/styles/macros_annotations.tex) -- reprise dans
% l'esprit du Dossier scientifique, à réutiliser/étendre à volonté.
\definecolor{colorOne}{HTML}{8C0000}
\definecolor{colorTwo}{HTML}{0F4C3A}
\definecolor{colorThree}{HTML}{1565C0}
\definecolor{colorFour}{HTML}{57659E}
\definecolor{colorFive}{HTML}{4A3628}
\definecolor{colorSix}{HTML}{2E7D5B}

% ============================================================
% TITRES PÉDAGOGIQUES AVEC MÉTADONNÉES
% ============================================================

\usepackage{etoolbox}

\usetikzlibrary{
  positioning,
  arrows.meta,
  shapes.geometric,
  calc
}


% ============================================================
% CLÉS DES MÉTADONNÉES
% ============================================================

\pgfkeys{
  /TitrePedagogique/.is family,
  /TitrePedagogique/indispensable/.initial=false,
  /TitrePedagogique/difficulte/.initial=,
  /TitrePedagogique/duree/.initial=,
  /TitrePedagogique/type/.initial=
}


% ============================================================
% PARAMÈTRES VISUELS DES MÉTADONNÉES
% ============================================================

\newcommand{\TitreMetaTailleType}{7.5pt}
\newcommand{\TitreMetaInterligneType}{9pt}

% Boîte et longueur utilisées pour MESURER la largeur réelle du bloc
% de métadonnées (flèche + étoiles + point + durée) et garantir que
% le trait gauche est toujours au moins aussi long que ce bloc (voir
% \TitrePedagogiqueAvecMeta). Déclarées une seule fois ici (pas dans
% la macro, sinon erreur "already defined" au 2e appel).
\newsavebox{\TitreMetaBox}
\newlength{\TitreBarGaucheNominale}
\newlength{\TitreBarGaucheLongueur}
% Boîte et longueur dédiées à la mesure de la largeur RÉELLE du
% titre (voir \TitrePedagogiqueAvecMeta) : on évite \sbox0/\wd0 (le
% registre de boîte global n°0, potentiellement réutilisé en
% coulisses par d'autres packages comme tcolorbox) au profit d'une
% boîte nommée. On stocke aussi le résultat dans un VRAI registre de
% longueur (\newlength) plutôt que dans une macro \def contenant
% \wd0 : le package calc (chargé par animate.sty) réécrit
% \addtolength/\setlength pour comprendre des expressions comme
% "0.5\linewidth", mais échoue à développer une macro qui contient
% elle-même une primitive \wd0 cachée -> "Illegal unit of measure"
% / "Package calc Error". Un \newlength classique, lui, est toujours
% géré nativement, sans ambiguïté.
\newsavebox{\TitreMesureTitre}
\newlength{\TitreLargeurEffLen}

% ============================================================
% LARGEUR ÉTROITE D'UN TITRE MULTI-LIGNES (remplace varwidth)
% ============================================================
%
% ⚠️ BUG DE FOND IDENTIFIÉ (et confirmé par isolation) : le package
% `anyfontsize` (chargé par templates/styles/monstyletikz.tex) patche
% une macro interne du noyau LaTeX (\tryif@simple, utilisée par la
% sélection de police) qui réutilise des registres temporaires
% globaux. Or `varwidth` a besoin de re-typographier le même
% paragraphe PLUSIEURS FOIS à des largeurs différentes pour trouver
% la plus étroite qui tienne -- donc rappelle \selectfont plusieurs
% fois -- et ces appels répétés, une fois patchés par anyfontsize,
% corrompent la recherche de varwidth, qui abandonne silencieusement
% et retombe sur la largeur maximale demandée (d'où le trait qui
% restait loin du titre quand celui-ci passait sur plusieurs lignes).
%
% \TitreLargeurEtroite ci-dessous fait le même travail que varwidth
% (trouver la largeur la plus étroite qui donne le même nombre de
% lignes qu'à la largeur maximale) mais SANS jamais rappeler
% \selectfont plus d'UNE fois -- donc immunisée contre ce conflit.
% Principe : la police est sélectionnée UNE SEULE fois (#3), puis on
% teste plusieurs largeurs candidates par dichotomie en comparant la
% HAUTEUR du paragraphe résultant à la hauteur de référence (largeur
% max) -- une ligne de plus fait sauter la hauteur d'un
% \baselineskip, donc "même hauteur" = "même nombre de lignes",
% sans avoir besoin de compter les lignes explicitement (ce qui,
% testé, s'est avéré bien plus fragile à manipuler au niveau TeX).
\newdimen\TitreDichoBasse
\newdimen\TitreDichoHaute
\newdimen\TitreDichoMilieu
\newdimen\TitreDichoHauteurRef
\newdimen\TitreDichoHauteurEssai
\newcount\TitreDichoIter

\newcommand{\TitreLargeurEtroite}[3]{%
  % #1 = largeur maximale (au-delà, retour à la ligne -- déjà connue
  %      comme donnant le nombre de lignes de référence)
  % #2 = texte du titre
  % #3 = code de sélection de police (ex: \fontsize{15}{17}\selectfont\bfseries)
  \begingroup
    #3\relax
    \setbox0=\vbox{\hsize=#1\relax \noindent #2\par}%
    \TitreDichoHauteurRef=\ht0 \advance\TitreDichoHauteurRef by \dp0
    \advance\TitreDichoHauteurRef by 0.5pt % petite marge de sécurité (arrondis)

    \TitreDichoBasse=0pt   % borne basse : peut donner TROP de lignes
    \TitreDichoHaute=#1\relax % borne haute : SÛRE (nb de lignes de référence)
    \TitreDichoIter=0
    \loop
      \TitreDichoMilieu=\TitreDichoBasse
      \advance\TitreDichoMilieu by \TitreDichoHaute
      \divide\TitreDichoMilieu by 2
      \setbox0=\vbox{\hsize=\TitreDichoMilieu\relax \noindent #2\par}%
      \TitreDichoHauteurEssai=\ht0 \advance\TitreDichoHauteurEssai by \dp0
      \ifdim\TitreDichoHauteurEssai>\TitreDichoHauteurRef
        \TitreDichoBasse=\TitreDichoMilieu % trop étroit (ligne en plus) -> remonter la borne basse
      \else
        \TitreDichoHaute=\TitreDichoMilieu % encore assez large -> peut resserrer encore
      \fi
      \advance\TitreDichoIter by 1
    \ifnum\TitreDichoIter<14 \repeat
    % ⚠️ BUG (encore un !) IDENTIFIÉ : \global\setlength{...}{...} ne
    % se propage PAS correctement quand \setlength est patché par le
    % package `calc` (chargé par templates/styles/packages.tex via
    % `animate`) -- confirmé par isolation. Affectation directe du
    % registre à la place (contourne \setlength entièrement) :
    \global\TitreLargeurEffLen=\TitreDichoHaute\relax
  \endgroup
}

% Boîte et longueur pour le cas "étoiles déplacées dans la barre
% DROITE" (voir plus bas : si le titre est trop long pour tenir sur
% une ligne, les étoiles quittent la barre gauche -- où restent
% flèche/point/durée -- pour se loger dans la barre droite, ce qui
% laisse le titre libre de passer sur plusieurs lignes).
\newsavebox{\TitreMetaBoxDroite}
\newlength{\TitreBarDroiteLongueur}
% Drapeaux (booléens TeX NATIFS, \newif) lus/positionnés par
% \TitrePedagogiqueAvecMeta et \TitreMetaBarreGauche.
%
% ⚠️ BUG CORRIGÉ : ces drapeaux étaient auparavant stockés comme des
% chaînes "true"/"false" et testés/positionnés via
% \ifstrequal{...}{true}{\def\Drapeau{false}}{...}. Ce motif s'est
% avéré NE PAS persister de façon fiable (la valeur revenait à son
% défaut peu après avoir été positionnée) -- c'est très probablement
% ce qui empêchait aussi bien la flèche "indispensable" que le
% déplacement des étoiles de fonctionner. Un \newif est un vrai
% conditionnel TeX (\iftrue/\iffalse en interne), sans ce genre
% d'ambiguïté : on l'utilise maintenant partout où on testait
% auparavant un drapeau "maison".
\newif\ifTitrePedTropLong
\newif\ifTitrePedEtoilesGauche
\TitrePedEtoilesGauchetrue % défaut : étoiles à gauche (cas normal)
\newif\ifTitrePedIndispensable

% Registres/booléens pour \TitreEtoilesDansBarrePath (mode normal
% 0-5 et mode binaire 6-31 -- voir cette macro pour le détail).
\newcount\TitreCountNiveau
\newcount\TitreCountQuot
\newif\ifTitreBinMode
\newif\ifTitreEtoileOui


% ⚠️ BUG CORRIGÉ (3e tentative) : mettre \ExplSyntaxOn / \ExplSyntaxOff
% À L'INTÉRIEUR du corps de \newcommand{...}{...} ne sert à RIEN :
% \newcommand tokénise (découpe en unités lexicales) tout son corps
% UNE SEULE FOIS, au moment de la DÉFINITION -- pas à chaque appel.
% Au moment de cette définition, les catcodes normaux étaient actifs
% (pas ceux d'expl3), donc "\tl_set:Nn" a été découpé en "\tl"
% (commande inconnue) + les caractères littéraux "_","s","e","t",...
% Un \ExplSyntaxOn PLACÉ DANS le corps arrive trop tard : le corps
% est déjà figé en tokens séparés, l'exécuter plus tard ne les
% "refusionne" pas. Il faut englober la DÉFINITION ELLE-MÊME (tout
% le \newcommand{...}{...}) entre \ExplSyntaxOn et \ExplSyntaxOff.
\ExplSyntaxOn

\tl_new:N \g_titrepedagogique_cles_tl

% Tampon réutilisé pour construire le texte de clés déjà 100%
% développé (voir \TitrePedagogiqueAppliquerCles ci-dessous).
\newcommand{\TitrePedagogiqueClesTmp}{}

\newcommand{\TitrePedagogiqueAppliquerCles}[1]{%
  \begingroup
    % \tl_trim_spaces:N retire l'espace/retour à la ligne que le
    % générateur .tex.j2 laisse en tête de liste (juste après le "["
    % du template) -- sinon pgfkeys échoue sur " indispensable" (avec
    % un espace) comme au tout premier essai.
    \tl_set:Nn \g_titrepedagogique_cles_tl {#1}
    \tl_trim_spaces:N \g_titrepedagogique_cles_tl
    % \xdef (edef + global) : on construit ICI, dans ce groupe, le
    % texte final "/TitrePedagogique/.cd,indispensable=True,..."
    % ENTIÈREMENT développé (plus aucune macro cachée dedans), et on
    % le rend global (\xdef) pour qu'il survive à la fermeture du
    % groupe juste en dessous.
    \xdef\TitrePedagogiqueClesTmp{/TitrePedagogique/.cd,\g_titrepedagogique_cles_tl}%
  \endgroup
  % \expandafter\pgfkeys\expandafter{\TitrePedagogiqueClesTmp} :
  % \TitrePedagogiqueClesTmp est une macro TeX "plate" déjà 100%
  % développée -- ce double \expandafter garantit que pgfkeys reçoit
  % directement les caractères bruts (avec de vraies virgules au
  % niveau top) AVANT même de commencer à lire son argument -- donc
  % son découpage par virgule fonctionne normalement, exactement
  % comme pour n'importe quel \pgfkeys{texte écrit à la main}.
  \expandafter\pgfkeys\expandafter{\TitrePedagogiqueClesTmp}%
}

\ExplSyntaxOff


% ============================================================
% ÉLÉMENTS "PATH" (à dessiner À L'INTÉRIEUR d'une scène TikZ déjà
% ouverte, sans créer leur propre \tikz{...}). Chaque élément est
% dessiné en partant de x=0 = son BORD GAUCHE, et est centré
% verticalement sur y=0. C'est le fait que TOUS les éléments soient
% tracés dans la MÊME scène, sur le MÊME y=0, qui garantit
% l'alignement parfait entre flèche / étoiles / point / durée / trait.
% ============================================================

% --- Flèche "Indispensable" ---------------------------------------
%
%   →
%
% #1 = Hauteur utile \HUtile (= hauteur totale de la flèche)
%
% ⚠️ Le facteur "1.5" ci-dessous (longueur de la flèche par rapport
% à sa hauteur) est repris tel quel dans \TitreMetaBarreGauche pour
% avancer le curseur \XCur après la flèche. Si tu changes ce 1.5
% ICI, pense à le changer aussi là-bas (sinon même bug que les
% étoiles : chevauchement avec l'élément suivant).
% ⚠️ BUG TIKZ CLASSIQUE (celui qui empêchait la flèche de s'afficher
% depuis le tout début, identifié et confirmé par isolation à la
% toute fin) : utiliser un PARAMÈTRE DE MACRO (#1) directement à
% l'intérieur d'une spécification de pointe de flèche TikZ
% "-{Latex[length=#1 pt, ...]}" est fragile -- TikZ analyse cette
% chaîne avec son propre tokenizer bas niveau, qui ne gère pas
% correctement un "#1" brut à cet endroit précis (contrairement à
% "\pgfmathsetmacro{...}{#1*1.5}" juste au-dessus, qui lui n'a
% jamais posé de problème). Le correctif standard : capturer #1
% dans une macro normale AVANT de l'utiliser dans la spec de flèche.
\newcommand{\TitreIndispensableDansBarrePath}[1]{%
  \edef\TitreArrowHUtile{#1}%
  \pgfmathsetmacro{\ArrowLen}{\TitreArrowHUtile*1.5}%
  \draw[
    -{Latex[length=\TitreArrowHUtile pt, width=\TitreArrowHUtile pt]},
    color=white,
    line width={\TitreArrowHUtile*0.4},
    line cap=butt %round rect
  ]
    (0,0) -- (\ArrowLen pt, 0);
}

% --- Étoiles de difficulté ------------------------------------------
%
%   ★ ★ ☆ ☆ ☆   (difficulté 1-5, mode normal)
%   ★ ☆ ★ ☆ ★   (difficulté 6-31, mode binaire -- voir plus bas)
%
% #1 = Niveau (0 à 31 -- 0 à 5 = mode normal, 6 à 31 = mode binaire)
% #2 = Couleur "creuse" (= couleur du titre / de la barre -- utilisée
%      pour les étoiles NON atteintes, qui se fondent dans la barre)
% #3 = Hauteur utile \HUtile (= taille "slot" de chaque étoile ;
%       la taille RÉELLEMENT dessinée est légèrement plus petite,
%       voir \TitreEtoileTailleFacteur ci-dessous)
%
% ⚠️ BUG CORRIGÉ : les étoiles "atteintes" (i<=niveau) étaient
% dessinées avec fill=#2 -- la MÊME couleur que le fond de la barre
% -- donc quasi invisibles dessus (seul un mince contour blanc les
% distinguait). Les étoiles "non atteintes" étaient en blanc plein,
% donc BEAUCOUP plus visibles -- l'inverse de ce qu'on veut : plus la
% difficulté est élevée, plus on doit voir de blanc plein, pas
% l'inverse. Les deux couleurs de remplissage sont donc échangées
% ci-dessous (\TitreEtoileRemplieCouleur / \TitreEtoileCreuseCouleur).
%
% 🎲 NOUVEAU -- mode binaire (difficulté 6 à 31) : pour une "vraie"
% difficulté au-delà de 5, les 5 étoiles représentent le NIVEAU EN
% BINAIRE sur 5 bits (poids 16-8-4-2-1, bit de poids fort à gauche),
% étoile pleine = bit à 1, étoile creuse = bit à 0. Permet d'aller
% jusqu'à 31 (= 11111 en binaire = 5 étoiles pleines).
%
% IMPORTANT : cette macro calcule et EXPOSE elle-même la largeur
% totale occupée (\TitreEtoilesLargeur), à partir du MÊME facteur
% d'espacement que celui utilisé pour le dessin. Comme ça, si tu
% changes \TitreEtoilesEspacementFacteur, la largeur suit
% automatiquement : plus jamais de désynchronisation avec ce qui
% vient après (la durée) -> c'est ce qui causait le chevauchement.
\newcommand{\TitreEtoilesDansBarrePath}[3]{%
  % Distance (en unités de \HUtile) entre les CENTRES de deux
  % étoiles consécutives. -> SEULE valeur à modifier pour changer
  % l'espacement entre étoiles.
  \def\TitreEtoilesEspacementFacteur{2.00}%
  % Taille RÉELLE de chaque étoile = ce facteur * \HUtile (facteur
  % légèrement < 1 pour compenser le fait que la forme "star" de
  % TikZ déborde un peu de sa boîte "minimum size" à cause de ses
  % pointes -> évite que les étoiles dépassent du trait).
  \def\TitreEtoileTailleFacteur{0.85}%

  % Niveau demandé, arrondi/borné à un entier entre 0 et 31 (garde-fou
  % : une valeur invalide ne doit jamais faire planter le document).
  \TitreCountNiveau=#1\relax
  \ifnum\TitreCountNiveau<0 \TitreCountNiveau=0 \fi
  \ifnum\TitreCountNiveau>31 \TitreCountNiveau=31 \fi
  % Mode binaire dès que le niveau dépasse 5 (la représentation
  % "N étoiles pleines sur 5" n'a plus de sens au-delà).
  \ifnum\TitreCountNiveau>5
    \TitreBinModetrue
  \else
    \TitreBinModefalse
  \fi

  \foreach \i/\TitrePoidsBit in {1/16,2/8,3/4,4/2,5/1}{%
    \pgfmathsetmacro{\XPos}{(\i-1)*(#3*\TitreEtoilesEspacementFacteur)}%
    \pgfmathsetmacro{\TailleEtoile}{#3*\TitreEtoileTailleFacteur}%
    % Détermine si CETTE étoile doit être "pleine" (\TitreEtoileOui)
    % ou "creuse", selon le mode normal ou binaire.
    \ifTitreBinMode
      \TitreCountQuot=\TitreCountNiveau
      \divide\TitreCountQuot by \TitrePoidsBit
      \ifodd\TitreCountQuot \TitreEtoileOuitrue \else \TitreEtoileOuifalse \fi
    \else
      \ifnum\i>\TitreCountNiveau \TitreEtoileOuifalse \else \TitreEtoileOuitrue \fi
    \fi
    \ifTitreEtoileOui
      % Étoile "atteinte" (pleine) : bien visible, blanc plein.
      \node[
        star, star points=5, star point ratio=.50, rounded corners=0.2pt,
        minimum size=\TailleEtoile pt, inner sep=0pt, outer sep=0pt,
        draw=white, fill=white, line width=0.4pt, rotate=180
      ] at (\XPos pt, 0) {};
    \else
      % Étoile "non atteinte" (creuse) : se fond dans la barre, ne
      % laisse qu'un mince contour blanc.
      \node[
        star, star points=5, star point ratio=.50, rounded corners=0.2pt,
        minimum size=\TailleEtoile pt, inner sep=0pt, outer sep=0pt,
        draw=white, fill=#2, line width=0.4pt, rotate=180
      ] at (\XPos pt, 0) {};
    \fi
  }%

  % Largeur totale du bloc = position de la dernière étoile (centre)
  % + sa propre taille "slot" (#3, pas la taille réduite, pour garder
  % une marge visuelle propre après la 5e étoile).
  \pgfmathsetmacro{\TitreEtoilesLargeurTmp}{4*(#3*\TitreEtoilesEspacementFacteur) + #3}%
  \xdef\TitreEtoilesLargeur{\TitreEtoilesLargeurTmp}%
}

% --- Séparateur (point) ---------------------------------------------
%
%   •
%
% #1 = Épaisseur de la barre (ex: 0.3cm)
% Diamètre du point = 20% de l'épaisseur de la barre -> rayon = 10%
\newcommand{\TitreSeparateurDansBarrePath}[1]{%
  \pgfmathsetmacro{\Rayon}{#1*0.1}%
  \fill[white] (\Rayon pt, 0) circle (\Rayon pt);
}


% ============================================================
% MÉTADONNÉES INCRUSTÉES DANS LA BARRE GAUCHE
% ============================================================
%
% Schéma global de la barre gauche, tout est sur le même axe y=0 :
%
%   trait ───────────────────────────────
%          →   ★ ★ ☆ ☆ ☆   •  25 min
%
\newcommand{\TitreMetaBarreGauche}[2]{%
  % #1 = Couleur de fond de la barre (utilisée pour les étoiles pleines)
  % #2 = Épaisseur de la barre (ex: 0.3cm)

  \edef\TitreTmpInd{\pgfkeysvalueof{/TitrePedagogique/indispensable}}%
  \edef\TitreTmpDiff{\pgfkeysvalueof{/TitrePedagogique/difficulte}}%
  \edef\TitreTmpDuree{\pgfkeysvalueof{/TitrePedagogique/duree}}%

  % --------------------------------------------------------------
  % Test "indispensable = vrai" ROBUSTE face à la casse.
  %
  % ⚠️ L'ancien code faisait :
  %     \lowercase{\edef\TitreTmpInd{\TitreTmpInd}}
  % ... ce qui ne fonctionne PAS : \lowercase ne lit que les
  % caractères LITTÉRALEMENT présents dans son argument, pas le
  % contenu caché à l'intérieur d'une macro. Comme Jinja/Python
  % écrit souvent "True"/"False" (majuscule) pour un booléen, la
  % comparaison à "true" échouait silencieusement -> la flèche ne
  % s'affichait jamais. On teste maintenant explicitement les
  % variantes usuelles, et on stocke le résultat dans un vrai
  % booléen TeX (\ifTitrePedIndispensable) plutôt que dans une
  % chaîne "true"/"false" re-testée ensuite via \ifstrequal -- ce
  % second motif s'est avéré ne pas persister de façon fiable.
  % --------------------------------------------------------------
  \TitrePedIndispensablefalse
  % ⚠️ BUG (enfin) IDENTIFIÉ : \ifstrequal{\Macro}{chaîne}{...}{...}
  % échoue silencieusement à exécuter une commande \newif comme code
  % vrai/faux dès qu'un des deux côtés comparés est une macro (même
  % après \edef) -- testé et reproduit en isolation. \ifdefstring
  % (etoolbox également, faite spécifiquement pour "comparer une
  % macro à une chaîne") n'a PAS ce problème.
  \ifdefstring{\TitreTmpInd}{true}{\TitrePedIndispensabletrue}{}%
  \ifdefstring{\TitreTmpInd}{True}{\TitrePedIndispensabletrue}{}%
  \ifdefstring{\TitreTmpInd}{1}{\TitrePedIndispensabletrue}{}%

  % Hauteur utile COMMUNE à la flèche, aux étoiles et à la durée.
  %
  % ⚠️ BUG CORRIGÉ : l'ancienne formule "épaisseur - 1mm" devient
  % NÉGATIVE dès que la barre est plus fine que 1mm -- c'est le cas
  % des niveaux 4 (0.12cm) et surtout 5 (0.08cm). Une longueur de
  % flèche négative fait planter en interne le calcul de la pointe
  % TikZ "Latex[...]" -> "Package PGF Math Error: 1/0.0". On passe
  % donc à une hauteur PROPORTIONNELLE (facteur de l'épaisseur),
  % qui reste TOUJOURS positive quelle que soit l'épaisseur.
  \pgfmathsetmacro{\HUtile}{#2*0.8}%

  \hbox{%
    \hspace{0.4em}% Marge gauche dans la barre
    % baseline=0pt : le y=0 de CETTE scène coïncide avec la ligne de
    % base courante -> c'est ce point qui sera ensuite réutilisé par
    % \TitrePedagogiqueAvecMeta pour aligner tout le monde sur y=0.
    \tikz[baseline=0pt]{%
      \pgfmathsetmacro{\XCur}{0}% curseur horizontal courant

      % --- 1. Flèche "Indispensable" ----------------------------
      %   →
      \ifTitrePedIndispensable
        \begin{scope}[xshift=\XCur pt]
          \TitreIndispensableDansBarrePath{\HUtile}%
        \end{scope}
        \pgfmathsetmacro{\XCur}{\XCur + 1.5*\HUtile + 0.5*\HUtile}%
      \fi

      % --- 2. Étoiles de difficulté ---------------------------------
      %   ★ ★ ☆ ☆ ☆
      % \ifTitrePedEtoilesGauche : positionné par
      % \TitrePedagogiqueAvecMeta. Si le titre est trop long pour
      % tenir sur une ligne, il vaut faux -- les étoiles sont alors
      % dessinées dans la barre DROITE à la place (voir
      % \TitreMetaBarreDroite), et on saute tout ce bloc ici (donc
      % le point+durée qui suit s'enchaîne directement après la
      % flèche, sans "trou" pour des étoiles absentes).
      \ifTitrePedEtoilesGauche
      \ifx\TitreTmpDiff\empty\else
        \begin{scope}[xshift=\XCur pt]
          \TitreEtoilesDansBarrePath{\TitreTmpDiff}{#1}{\HUtile}%
        \end{scope}
        % \TitreEtoilesLargeur vient d'être calculée PAR
        % \TitreEtoilesDansBarrePath elle-même (voir plus haut) :
        % toujours synchronisée avec l'espacement réellement utilisé.
        % Espace entre la fin du bloc étoiles et le point de séparation,
        % en unités de \HUtile. -> C'est ICI qu'il faut augmenter la
        % valeur si tu veux encore plus d'écart avant le point/la durée.
        \def\TitreEspaceEtoilesPoint{1.2}%
        \pgfmathsetmacro{\XCur}{\XCur + \TitreEtoilesLargeur + \TitreEspaceEtoilesPoint*\HUtile}%
      \fi
      \fi

      % --- 3. Séparateur (point) puis durée -----------------------------
      %   •  25 min
      \ifx\TitreTmpDuree\empty\else
        \begin{scope}[xshift=\XCur pt]
          \TitreSeparateurDansBarrePath{#2}%
        \end{scope}
        \pgfmathsetmacro{\DiamPoint}{#2*0.2}%
        \pgfmathsetmacro{\XCur}{\XCur + \DiamPoint + 0.35*\HUtile}%
        \node[
          anchor=west, inner sep=0pt, outer sep=0pt,
          font={\fontsize{\HUtile pt}{\HUtile pt}\selectfont\bfseries},
          text=white
        ] at (\XCur pt, 0) {\TitreTmpDuree};%
      \fi
    }%
    \hspace{0.2em}%
  }%
}


% ============================================================
% MÉTADONNÉES INCRUSTÉES DANS LA BARRE DROITE (étoiles seules)
% ============================================================
%
% N'est utilisée QUE quand le titre est trop long pour tenir sur une
% ligne (voir \TitrePedagogiqueAvecMeta) : les étoiles quittent alors
% la barre gauche pour se loger ici, dans la barre droite, ce qui
% libère de la place à gauche et laisse le titre passer sur
% plusieurs lignes sans être davantage comprimé.
%
%   trait ─────────
%          ★ ★ ☆ ☆ ☆
%
\newcommand{\TitreMetaBarreDroite}[2]{%
  % #1 = Couleur de fond de la barre (utilisée pour les étoiles pleines)
  % #2 = Épaisseur de la barre (ex: 0.3cm)
  \edef\TitreTmpDiffDroite{\pgfkeysvalueof{/TitrePedagogique/difficulte}}%
  \pgfmathsetmacro{\HUtile}{#2*0.8}%
  \hbox{%
    \hspace{0.2em}%
    \tikz[baseline=0pt]{%
      \ifx\TitreTmpDiffDroite\empty\else
        \TitreEtoilesDansBarrePath{\TitreTmpDiffDroite}{#1}{\HUtile}%
      \fi
    }%
    \hspace{0.4em}%
  }%
}


% ============================================================
% TYPE D'EXERCICE
% ============================================================

\newcommand{\TitreMetaType}[1]{%
  \edef\TitreTmpType{\pgfkeysvalueof{/TitrePedagogique/type}}%
  \ifx\TitreTmpType\empty\else
    \par\vspace{0.15ex}%
    {\centering
      \fontsize{\TitreMetaTailleType}{\TitreMetaInterligneType}\selectfont
      \itshape
      \textcolor{#1}{\TitreTmpType}%
      \par
    }%
  \fi
}


% ============================================================
% TITRE PÉDAGOGIQUE GÉNÉRIQUE
% ============================================================

\newcommand{\TitrePedagogiqueAvecMeta}[7]{%
  % #1 = couleur
  % #2 = police du titre (ex: \fontsize{15}{17}\selectfont\bfseries)
  % #3 = largeur MAXIMALE du titre (au-delà, retour à la ligne)
  % #4 = épaisseur des barres
  % #5 = distance entre barre et titre
  % #6 = espace vertical après le titre
  % #7 = texte du titre

  % --------------------------------------------------------------
  % Le titre a-t-il BESOIN de passer à la ligne ? (comparaison de sa
  % largeur naturelle, SUR UNE SEULE LIGNE, sans limite, à #3). Sert
  % à décider où vont les étoiles (voir \TitreEtoilesGaucheOui juste
  % en dessous) : \hbox n'ayant pas de limite de largeur, c'est la
  % seule façon fiable de savoir si le texte "voudrait" tenir sur une
  % ligne ou pas, indépendamment de la largeur repliée qu'on calcule
  % juste après avec varwidth.
  % --------------------------------------------------------------
  \sbox{\TitreMesureTitre}{\hbox{#2#7}}%
  \ifdim\wd\TitreMesureTitre>#3
    \TitrePedTropLongtrue
  \else
    \TitrePedTropLongfalse
  \fi

  % Étoiles à gauche (cas normal) sauf si le titre est trop long :
  % dans ce cas elles migrent vers la barre DROITE (voir
  % \TitreMetaBarreDroite plus bas) pour ne pas comprimer davantage
  % un titre qui a de toute façon déjà besoin de plusieurs lignes.
  \ifTitrePedTropLong
    \TitrePedEtoilesGauchefalse
  \else
    \TitrePedEtoilesGauchetrue
  \fi

  % --------------------------------------------------------------
  % Largeur RÉELLE du titre.
  %
  % ⚠️ Ne dépend PLUS de varwidth (voir \TitreLargeurEtroite,
  % déclarée plus haut, et son commentaire pour le pourquoi complet
  % du conflit avec anyfontsize). Si le titre tient sur une ligne :
  % on réutilise directement sa largeur naturelle, déjà mesurée
  % ci-dessus (aucun calcul supplémentaire). Sinon : dichotomie
  % maison, immunisée contre le conflit anyfontsize/varwidth.
  % --------------------------------------------------------------
  \ifTitrePedTropLong
    \TitreLargeurEtroite{#3}{#7}{#2}%
  \else
    \setlength{\TitreLargeurEffLen}{\wd\TitreMesureTitre}%
  \fi

  % --------------------------------------------------------------
  % Mesure du bloc de métadonnées, ELLE AUSSI avant l'ouverture du
  % tikzpicture principal (comme le titre juste au-dessus).
  %
  % ⚠️ BUG CORRIGÉ : \TitreMetaBarreGauche crée SA PROPRE scène TikZ
  % interne (\tikz[baseline=0pt]{...}). La mesurer via \sbox AU
  % MILIEU d'un tikzpicture déjà ouvert (comme c'était fait avant,
  % dans l'étape 3 plus bas) empile une scène TikZ dans une boîte
  % PENDANT qu'une autre est activement en train de dessiner : ça
  % perturbe la pile de polices interne de pgf et la fait boucler à
  % l'infini -> "TeX capacity exceeded" sur \pgf@selectfontorig. En
  % la mesurant ICI, hors de tout tikzpicture actif (exactement
  % comme pour le titre), on évite ce conflit.
  \sbox{\TitreMetaBox}{\TitreMetaBarreGauche{#1}{#4}}%

  % Mesure de la boîte "étoiles pour la barre droite", SEULEMENT si
  % elles y ont été déplacées (titre trop long) ET qu'il y a
  % effectivement une difficulté à afficher (sinon rien à mesurer).
  \edef\TitreTmpDiffCheck{\pgfkeysvalueof{/TitrePedagogique/difficulte}}%
  \ifTitrePedTropLong
    \ifx\TitreTmpDiffCheck\empty\else
      \sbox{\TitreMetaBoxDroite}{\TitreMetaBarreDroite{#1}{#4}}%
    \fi
  \fi

  \noindent
  \begin{tikzpicture}[baseline=(title.base)]

    % --------------------------------------------------------
    % 1. TITRE PRINCIPAL
    % --------------------------------------------------------
    %        |<-#5->[ TITRE ]<-#5->|
    %
    % Ancre par défaut = "center" : title est donc centré
    % verticalement sur y=0, et title.west/east collent au texte
    % réel grâce à \TitreLargeurEffLen.
    \node[
      inner sep=0pt,
      outer sep=0pt,
      align=center,
      text width=\TitreLargeurEffLen,
      font=#2,
      text=#1
    ] (title) at (0.5\linewidth,0) {#7};


    % --------------------------------------------------------
    % 2. TYPE D'EXERCICE (Sous le titre)
    % --------------------------------------------------------
    \node[
      below=0.25ex of title.south,
      inner sep=0pt,
      outer sep=0pt,
      text width=#3,
      align=center
    ] (type) {%
      \TitreMetaType{#1}%
    };


    % --------------------------------------------------------
    % 3. BARRE GAUCHE (trait + métadonnées incrustées)
    % --------------------------------------------------------
    %
    %  trait ───────────────────────  [ TITRE ]
    %   →  ★★☆☆☆   • 25 min     <-->
    %                             #5
    %
    % IMPORTANT : le trait doit être AU MOINS aussi long que le bloc
    % de métadonnées, sinon la fin du bloc (flèche, étoiles, point,
    % durée -> tout est dessiné en BLANC) déborde sur le fond
    % blanc/transparent de la page et devient invisible. On mesure
    % donc la largeur RÉELLE du bloc (\TitreMetaBox) et on prend le
    % maximum entre cette largeur et la longueur "nominale" (celle
    % qui amène pile jusqu'au titre).

    % Largeur "nominale" du trait gauche = distance entre le bord
    % gauche (x=0) et le bord gauche du titre, moins l'espace #5.
    % Le titre est centré sur 0.5\linewidth avec une largeur
    % \TitreLargeurEffLen -> son bord gauche est à
    % 0.5\linewidth - 0.5\TitreLargeurEffLen.
    \setlength{\TitreBarGaucheNominale}{0.5\linewidth}%
    \addtolength{\TitreBarGaucheNominale}{-0.5\TitreLargeurEffLen}%
    \addtolength{\TitreBarGaucheNominale}{-#5}%

    % Largeur réelle du bloc de métadonnées : déjà mesurée plus haut,
    % AVANT l'ouverture de ce tikzpicture (voir \TitreMetaBox) --
    % on ne la remesure pas ici, on compare juste au trait nominal.
    \ifdim\wd\TitreMetaBox>\TitreBarGaucheNominale
      \setlength{\TitreBarGaucheLongueur}{\wd\TitreMetaBox}%
    \else
      \setlength{\TitreBarGaucheLongueur}{\TitreBarGaucheNominale}%
    \fi
    \pgfmathsetmacro{\BarGaucheLongueurPt}{\TitreBarGaucheLongueur/1pt}%

    \draw[
      color=#1,
      line width=#4,
      line cap=rect
    ] (0,0) -- (\BarGaucheLongueurPt pt, 0);

    % anchor=base west : force le y=0 INTERNE de \TitreMetaBarreGauche
    % (fixé par \tikz[baseline=0pt] dans cette macro) à coïncider
    % EXACTEMENT avec le y=0 de cette scène-ci -> même axe que le
    % trait et que le centre du titre, quel que soit le contenu.
    % On réutilise directement la boîte déjà mesurée (\usebox), pas
    % besoin de redessiner le contenu une seconde fois.
    \node[
      anchor=base west,
      inner sep=0pt,
      outer sep=0pt
    ] at (0,0) {\usebox{\TitreMetaBox}};


    % --------------------------------------------------------
    % 4. BARRE DROITE (trait plein, ou étoiles si déplacées ici)
    % --------------------------------------------------------
    %
    %  [ TITRE ]      trait ──────────────────
    %            <-->
    %             #5
    %
    %  -- ou, si le titre est trop long (étoiles déplacées) --
    %
    %  [ TITRE  ]     trait ─────
    %  (2 lignes) <-->  ★ ★ ☆ ☆ ☆
    \path let \p1 = (title.east) in
      coordinate (barRightStart) at (\x1 + #5, 0);

    \ifTitrePedTropLong
      \ifx\TitreTmpDiffCheck\empty
        % titre trop long mais rien à afficher (pas de difficulté) :
        % trait droit normal, rien de spécial à faire.
        \draw[color=#1, line width=#4, line cap=rect] (barRightStart) -- (\linewidth,0);
      \else
        % Étoiles déplacées dans la barre droite : même logique que la
        % barre gauche (trait au moins aussi long que le contenu réel,
        % voir plus haut) -- comme le titre est centré, la longueur
        % "nominale" du côté droit est la MÊME que \TitreBarGaucheNominale
        % (symétrie du centrage), pas besoin de la recalculer.
        \ifdim\wd\TitreMetaBoxDroite>\TitreBarGaucheNominale
          \setlength{\TitreBarDroiteLongueur}{\wd\TitreMetaBoxDroite}%
        \else
          \setlength{\TitreBarDroiteLongueur}{\TitreBarGaucheNominale}%
        \fi
        \pgfmathsetmacro{\BarDroiteLongueurPt}{\TitreBarDroiteLongueur/1pt}%

        \draw[color=#1, line width=#4, line cap=rect]
          (barRightStart) -- ($(barRightStart)+(\BarDroiteLongueurPt pt,0)$);

        \node[anchor=base west, inner sep=0pt, outer sep=0pt]
          at (barRightStart) {\usebox{\TitreMetaBoxDroite}};
      \fi
    \else
      \draw[
        color=#1,
        line width=#4,
        line cap=rect
      ] (barRightStart) -- (\linewidth,0);
    \fi%

  \end{tikzpicture}

  \par
  \vspace{#6}
}


% ============================================================
% COMMANDES NIVEAUX 1 À 5
% ============================================================

\NewDocumentCommand{\TitreNiveauUn}{O{}m}{%
  \pgfkeys{/TitrePedagogique/indispensable=false, /TitrePedagogique/difficulte=, /TitrePedagogique/duree=, /TitrePedagogique/type=}%
  \TitrePedagogiqueAppliquerCles{#1}%
  \TitrePedagogiqueAvecMeta{colorStitre}{\fontsize{15}{17}\selectfont\bfseries}{0.65\linewidth}{0.3cm}{0.4cm}{0.35cm}{#2}%
}

\NewDocumentCommand{\TitreNiveauDeux}{O{}m}{%
  \pgfkeys{/TitrePedagogique/indispensable=false, /TitrePedagogique/difficulte=, /TitrePedagogique/duree=, /TitrePedagogique/type=}%
  \TitrePedagogiqueAppliquerCles{#1}%
  \TitrePedagogiqueAvecMeta{colorSStitre}{\fontsize{12}{14}\selectfont\bfseries}{0.65\linewidth}{0.22cm}{0.3cm}{0.30cm}{#2}%
}

\NewDocumentCommand{\TitreNiveauTrois}{O{}m}{%
  \pgfkeys{/TitrePedagogique/indispensable=false, /TitrePedagogique/difficulte=, /TitrePedagogique/duree=, /TitrePedagogique/type=}%
  \TitrePedagogiqueAppliquerCles{#1}%
  \TitrePedagogiqueAvecMeta{colorSSStitre}{\fontsize{10}{12}\selectfont\bfseries}{0.65\linewidth}{0.16cm}{0.2cm}{0.25cm}{#2}%
}

\NewDocumentCommand{\TitreNiveauQuatre}{O{}m}{%
  \pgfkeys{/TitrePedagogique/indispensable=false, /TitrePedagogique/difficulte=, /TitrePedagogique/duree=, /TitrePedagogique/type=}%
  \TitrePedagogiqueAppliquerCles{#1}%
  \TitrePedagogiqueAvecMeta{colorSSStitre}{\fontsize{10}{12}\selectfont\itshape}{0.65\linewidth}{0.12cm}{0.15cm}{0.20cm}{#2}%
}

\NewDocumentCommand{\TitreNiveauCinq}{O{}m}{%
  \pgfkeys{/TitrePedagogique/indispensable=false, /TitrePedagogique/difficulte=, /TitrePedagogique/duree=, /TitrePedagogique/type=}%
  \TitrePedagogiqueAppliquerCles{#1}%
  \TitrePedagogiqueAvecMeta{grayBrume}{\fontsize{9}{11}\selectfont\itshape}{0.65\linewidth}{0.08cm}{0.10cm}{0.15cm}{#2}%
}





% ----------------------------------------------------------

\newcommand{\TitreNiveauUnOld}[1]{%
  \noindent\begin{tikzpicture}
    \node[inner sep=0pt, align=center, font=\fontsize{15}{17}\selectfont\bfseries,
      color=colorStitre] (title) at (0.5\linewidth,0) {
      \begin{varwidth}{0.85\linewidth}\centering #1\end{varwidth}
    };
    \draw[colorStitre, line width=0.3cm] (0,0) -- ([xshift=-0.5cm]title.west);
    \draw[colorStitre, line width=0.3cm] ([xshift=0.5cm]title.east) -- (\linewidth,0);
  \end{tikzpicture}
  \par\vspace{0.35cm}
}

\newcommand{\TitreNiveauDeuxOld}[1]{%
  \noindent\begin{tikzpicture}
    \node[inner sep=0pt, align=center, font=\fontsize{12}{14}\selectfont\bfseries,
      color=colorSStitre] (title) at (0.5\linewidth,0) {
      \begin{varwidth}{0.85\linewidth}\centering #1\end{varwidth}
    };
    \draw[colorSStitre, line width=0.15cm] (0,0) -- ([xshift=-0.3cm]title.west);
    \draw[colorSStitre, line width=0.15cm] ([xshift=0.3cm]title.east) -- (\linewidth,0);
  \end{tikzpicture}
  \par\vspace{0.3cm}
}

\newcommand{\TitreNiveauTroisOld}[1]{%
  \noindent\begin{tikzpicture}
    \node[inner sep=0pt, align=center, font=\fontsize{10}{12}\selectfont\bfseries,
      color=colorSSStitre] (title) at (0.5\linewidth,0) {
      \begin{varwidth}{0.85\linewidth}\centering #1\end{varwidth}
    };
    \draw[colorSSStitre, line width=0.10cm] (0,0) -- ([xshift=-0.2cm]title.west);
    \draw[colorSSStitre, line width=0.10cm] ([xshift=0.2cm]title.east) -- (\linewidth,0);
  \end{tikzpicture}
  \par\vspace{0.25cm}
}

% Niveaux 4 et 5 : pas dans le fichier d'origine -- extrapolés dans le
% même esprit (barre de plus en plus fine, texte de plus en plus léger).
\newcommand{\TitreNiveauQuatreOld}[1]{%
  \noindent{\color{colorSSStitre}\rule{0.4cm}{0.06cm}}\ %
  \textit{\fontsize{10}{12}\selectfont #1}
  \par\vspace{0.2cm}
}

\newcommand{\TitreNiveauCinqOld}[1]{%
  \noindent{\color{grayBrume}\rule{0.25cm}{0.04cm}}\ %
  \textit{\fontsize{9}{11}\selectfont\color{grayBrume} #1}
  \par\vspace{0.15cm}
}

% ------------------------------------------------------------------
% \EncadreUn{titre vertical}{couleur}{contenu} et
% \EncadreDeux{titre vertical (centre)}{petit label (bas)}{couleur}{contenu}
% -- repris tels quels du Dossier scientifique. Disponibles pour tout
% usage libre dans les champs `enonce`/`correction`/`solution`/... du YAML.
% ------------------------------------------------------------------
\newcommand{\EncadreUn}[3]{%
  \begin{tcolorbox}[
      enhanced, notitle, colback=white, colbacklower=white,
      frame hidden, boxrule=0pt, sharp corners,
      borderline west={4pt}{0pt}{#2}, left=6mm, fontupper=\sffamily,
      overlay={
          \fill[#2] (frame.south west) rectangle ([xshift=4pt]frame.north west);
          \node[rotate=90, anchor=center, font=\sffamily\bfseries\small, text=white]
            at ([xshift=2pt]frame.west |- frame.center) {#1};
        }
    ]
    #3
  \end{tcolorbox}
}

\newcommand{\EncadreDeux}[4]{%
  \begin{tcolorbox}[
      enhanced, notitle, colback=white, colbacklower=white,
      frame hidden, boxrule=0pt, sharp corners,
      borderline west={4pt}{0pt}{#3}, left=7mm, fontupper=\sffamily,
      overlay={
          \fill[#3] (frame.south west) rectangle ([xshift=4pt]frame.north west);
          \node[rotate=90, anchor=south, font=\sffamily\bfseries, text=#3]
            at ([xshift=2pt]frame.west) {#2};
          \node[rotate=90, anchor=center, font=\sffamily\bfseries\small, text=white]
            at ([xshift=2pt]frame.west |- frame.center) {#1};
        }
    ]
    #4
  \end{tcolorbox}
}

% ============================================================
% Environnement générique
% ============================================================

\NewDocumentEnvironment{encadrestyle}{mmm}
{%
  \begin{tcolorbox}[
    %float, % <-- Rend la tcolorbox elle-même flottante
    enhanced,
    notitle,
    colback=#3!5,
    colbacklower=white,
    frame hidden,
    boxrule=0pt,
    sharp corners,
    borderline west={4pt}{0pt}{#3},
    left=7mm,
    fontupper=\sffamily,
    overlay={%
      % Bande verticale
      \fill[#3]
        (frame.south west)
        rectangle
        ([xshift=4pt]frame.north west);

      % Étiquette
      \node[
        rotate=90,
        anchor=south,
        font=\sffamily\bfseries\fontsize{5pt}{1pt}\selectfont\bfseries,
        text=#3,
        inner sep=0pt
      ] at ([xshift=-3pt]frame.west)
      {#2};

      % Titre
      \node[
        rotate=90,
        anchor=center,
        font=\sffamily\bfseries\small\fontsize{5pt}{10pt}\selectfont\bfseries,
        text=white,
        inner sep=0pt
      ] at ([xshift=2pt]frame.west |- frame.center)
      {#1};
    }
  ]
}
{%
  \end{tcolorbox}%
}


% ============================================================
% Environnements pédagogiques
% ============================================================


\NewDocumentEnvironment{prerequis}{o}
{%
  \IfNoValueTF{#1}
    {\begin{encadrestyle}{Prérequis}{}{colorOchre}}
    {\begin{encadrestyle}{Prérequis}{#1}{colorOchre}}
}
{%
  \end{encadrestyle}%
}

\NewDocumentEnvironment{objectifspedagogiquesetudiants}{o}
{%
  \IfNoValueTF{#1}
    {\begin{encadrestyle}{Objectifs pédagogiques (étudiants)}{}{colorObjectifsPedagogiquesEtudiants}}
    {\begin{encadrestyle}{Objectifs pédagogiques (étudiants)}{#1}{colorObjectifsPedagogiquesEtudiants}}
}
{%
  \end{encadrestyle}%
}

\NewDocumentEnvironment{objectifspedagogiquesenseignants}{o}
{%
  \IfNoValueTF{#1}
    {\begin{encadrestyle}{Objectifs pédagogiques (enseignants)}{}{colorObjectifsPedagogiquesEnseignants}}
    {\begin{encadrestyle}{Objectifs pédagogiques (enseignants)}{#1}{colorObjectifsPedagogiquesEnseignants}}
}
{%
  \end{encadrestyle}%
}

\NewDocumentEnvironment{pointcours}{o}
{%
  \IfNoValueTF{#1}
    {\begin{encadrestyle}{Point de cours}{}{colorPointCours}}
    {\begin{encadrestyle}{Point de cours}{#1}{colorPointCours}}
}
{%
  \end{encadrestyle}%
}

\NewDocumentEnvironment{coupdepouce}{o}
{%
  \IfNoValueTF{#1}
    {\begin{encadrestyle}{Coup de pouce}{}{colorOchre}}
    {\begin{encadrestyle}{Coup de pouce}{#1}{colorOchre}}
}
{%
  \end{encadrestyle}%
}


\NewDocumentEnvironment{correction}{o}
{%
  \IfNoValueTF{#1}
    {\begin{encadrestyle}{Correction}{}{greenSapin}}
    {\begin{encadrestyle}{Correction}{#1}{greenSapin}}
}
{%
  \end{encadrestyle}%
}


\NewDocumentEnvironment{solutiondetaillee}{o}
{%
  \IfNoValueTF{#1}
    {\begin{encadrestyle}{Solution détaillée}{}{grayPierre}}
    {\begin{encadrestyle}{Solution détaillée}{#1}{grayPierre}}
}
{%
  \end{encadrestyle}%
}


\NewDocumentEnvironment{remarque}{o}
{%
  \IfNoValueTF{#1}
    {\begin{encadrestyle}{Remarque}{}{headerblue}}
    {\begin{encadrestyle}{Remarque}{#1}{headerblue}}
}
{%
  \end{encadrestyle}%
}

\NewDocumentEnvironment{remarquespedagogiquesecueilsclassiques}{o}
{%
  \IfNoValueTF{#1}
    {\begin{encadrestyle}{Remarques pédagogiques (écueils classiques)}{}{colorRemarquesPedagogiquesEcueilsClassiques}}
    {\begin{encadrestyle}{Remarques pédagogiques (écueils classiques)}{#1}{colorRemarquesPedagogiquesEcueilsClassiques}}
}
{%
  \end{encadrestyle}%
}

\NewDocumentEnvironment{remarquespedagogiquesconseilsanimation}{o}
{%
  \IfNoValueTF{#1}
    {\begin{encadrestyle}{Remarques pédagogiques (conseils d'animation)}{}{colorRemarquesPedagogiquesConseilsAnimation}}
    {\begin{encadrestyle}{Remarques pédagogiques (conseils d'animation)}{#1}{colorRemarquesPedagogiquesConseilsAnimation}}
}
{%
  \end{encadrestyle}%
}

\NewDocumentEnvironment{bilan}{o}
{%
  \IfNoValueTF{#1}
    {\begin{encadrestyle}{Bilan}{}{colorBilan}}
    {\begin{encadrestyle}{Bilan}{#1}{colorBilan}}
}
{%
  \end{encadrestyle}%
}


\NewDocumentEnvironment{allerplusloin}{o}
{%
  \IfNoValueTF{#1}
    {\begin{encadrestyle}{Aller plus loin}{}{Prune}}
    {\begin{encadrestyle}{Aller plus loin}{#1}{Prune}}
}
{%
  \end{encadrestyle}%
}


% % ============================================================
% % TITRES PÉDAGOGIQUES AVEC MÉTADONNÉES
% % ============================================================

% \usepackage{etoolbox}

% \usetikzlibrary{
%   positioning,
%   arrows.meta,
%   shapes.geometric
% }


% % ============================================================
% % CLÉS DES MÉTADONNÉES
% % ============================================================

% % IMPORTANT :
% % On donne ici le chemin COMPLET à chaque clé.
% %
% % Cela évite toute ambiguïté avec :
% %
% %   indispensable/.initial
% %
% % et garantit que les clés appartiennent bien à la famille
% % /TitrePedagogique.

% \pgfkeys{
%   /TitrePedagogique/.is family,
%   /TitrePedagogique/indispensable/.initial=false,
%   /TitrePedagogique/difficulte/.initial=,
%   /TitrePedagogique/duree/.initial=,
%   /TitrePedagogique/type/.initial=
% }


% % ============================================================
% % PARAMÈTRES VISUELS DES MÉTADONNÉES
% % ============================================================

% % Ces paramètres sont modifiés automatiquement par
% % \TitrePedagogiqueAvecMeta.
% %
% % Ils permettent de régler séparément :
% %
% %   - taille des étoiles
% %   - taille de la flèche
% %   - taille de la durée
% %   - taille du type
% %   - espacement entre les éléments

% \newcommand{\TitreMetaTailleEtoile}{0.78em}
% \newcommand{\TitreMetaEchelleEtoile}{0.78}
% \newcommand{\TitreMetaTailleDuree}{8pt}
% \newcommand{\TitreMetaInterligneDuree}{9.5pt}
% \newcommand{\TitreMetaTailleType}{7.5pt}
% \newcommand{\TitreMetaInterligneType}{9pt}

% \newcommand{\TitreMetaEspaceFlecheEtoiles}{0.65em}
% \newcommand{\TitreMetaEspaceEtoilesDuree}{0.65em}
% \newcommand{\TitreMetaEspaceAvantFleche}{0.15em}
% \newcommand{\TitreMetaEspaceApresDuree}{0.15em}


% % ============================================================
% % ÉTOILES DE DIFFICULTÉ
% % ============================================================

% % #1 = niveau de difficulté (1 à 5)
% % #2 = couleur
% %
% % Les étoiles sont volontairement un peu espacées afin de ne
% % pas donner l'impression d'un seul bloc.
% %
% % Elles utilisent la même couleur que la barre et le titre.

% \newcommand{\TitreEtoiles}[2]{%
%   \tikz[
%     baseline=-0.35ex,
%     x=\TitreMetaEchelleEtoile em,
%     y=\TitreMetaEchelleEtoile em,
%   ]{%
%     \foreach \i in {1,...,5}{%
%       \ifnum\i>#1
%         % Étoile vide
%         \node[
%           star,
%           star points=5,
%           star point ratio=.50,
%           corner radius=0.5pt,
%           rounded corners=1.5pt, % Ajuste l'arrondi des coins (ex: 1.5pt)
%           minimum size=\TitreMetaTailleEtoile,
%           inner sep=0.05pt,
%           outer sep=0pt,
%           draw=#2,
%           fill=white,
%           line width=1.2pt,      % Ajuste l'épaisseur du trait (ex: 1.2pt)
%           rotate=180
%         ] at ({(\i-1)*2.2},0) {};
%       \else
%         % Étoile pleine
%         \node[
%           star,
%           star points=5,
%           star point ratio=.50,
%           corner radius=1.5pt,
%           rounded corners=1.5pt, % Ajuste l'arrondi des coins (ex: 1.5pt)
%           minimum size=\TitreMetaTailleEtoile,
%           inner sep=0.05pt,
%           outer sep=0pt,
%           draw=#2,
%           fill=#2,
%           line width=1.2pt,      % Ajuste l'épaisseur du trait (ex: 1.2pt)
%           rotate=180
%         ] at ({(\i-1)*2.2},0) {};
%       \fi
%     }%
%   }%
% }


% % ============================================================
% % FLÈCHE "INDISPENSABLE"
% % ============================================================

% % #1 = couleur
% %
% % Flèche volontairement petite et fine.
% % Elle est entièrement dessinée en TikZ.

% \newcommand{\TitreIndispensable}[1]{%
%   \tikz[
%     baseline=-0.35ex
%   ]{%
%     \draw[
%       -{Latex[
%         length=1.55cm,
%         width=1.15cm
%       ]
%       },
%       color=#1,
%       line width=0.65cm,
%     ]
%       (0,0) -- (0.95em,0);
%   }%
% }


% % ============================================================
% % SÉPARATEUR
% % ============================================================

% % Petit point entre les étoiles et la durée.

% \newcommand{\TitreSeparateur}[1]{%
%   \tikz[
%     baseline=-0.35ex
%   ]{%
%     \fill[
%       #1
%     ]
%       (0,0)
%       circle (0.20ex);
%   }%
% }


% % ============================================================
% % MÉTADONNÉES SUR LA LIGNE DU TITRE
% % ============================================================

% % Cette macro ne connaît PAS la couleur.
% % La couleur lui est transmise explicitement par
% % \TitrePedagogiqueAvecMeta.
% %
% % Ordre :
% %
% %   →   ★★☆☆☆   • 25 min
% %
% % Tous les éléments utilisent la même couleur que le titre.

% \newcommand{\TitreMetaInline}[1]{%

%   % ----------------------------------------------------------
%   % Récupération des valeurs
%   % ----------------------------------------------------------

%   \edef\TitreTmpInd{%
%     \pgfkeysvalueof{/TitrePedagogique/indispensable}%
%   }%

%   \edef\TitreTmpDiff{%
%     \pgfkeysvalueof{/TitrePedagogique/difficulte}%
%   }%

%   \edef\TitreTmpDuree{%
%     \pgfkeysvalueof{/TitrePedagogique/duree}%
%   }%


%   % ----------------------------------------------------------
%   % Normalisation de True / TRUE / true
%   % ----------------------------------------------------------

%   \lowercase{\edef\TitreTmpInd{\TitreTmpInd}}%


%   % ----------------------------------------------------------
%   % INDICATEUR "INDISPENSABLE"
%   % ----------------------------------------------------------

%   \ifstrequal{\TitreTmpInd}{true}{%
%     %
%     % Petite flèche
%     %
%     \TitreIndispensable{#1}%
%     \hspace{\TitreMetaEspaceFlecheEtoiles}%
%   }{}%


%   % ----------------------------------------------------------
%   % DIFFICULTÉ
%   % ----------------------------------------------------------

%   \ifx\TitreTmpDiff\empty
%     % Pas de difficulté : rien
%   \else
%     %
%     % Étoiles
%     %
%     \TitreEtoiles{\TitreTmpDiff}{#1}%
%     %
%     \hspace{\TitreMetaEspaceEtoilesDuree}%
%   \fi


%   % ----------------------------------------------------------
%   % DURÉE
%   % ----------------------------------------------------------

%   \ifx\TitreTmpDuree\empty
%     % Pas de durée : rien
%   \else
%     %
%     % Petit point
%     %
%     \TitreSeparateur{#1}%
%     %
%     \hspace{0.40em}%
%     %
%     % Durée
%     %
%     {\fontsize{\TitreMetaTailleDuree}
%              {\TitreMetaInterligneDuree}
%      \selectfont
%      \textcolor{#1}{\TitreTmpDuree}%
%     }%
%     %
%     \hspace{\TitreMetaEspaceApresDuree}%
%   \fi
% }


% % ============================================================
% % TYPE D'EXERCICE
% % ============================================================

% % Le type est placé sous le titre.
% %
% % Il utilise la même couleur que la barre et le titre,
% % mais avec une taille plus petite et une graisse légère.

% \newcommand{\TitreMetaType}[1]{%

%   \edef\TitreTmpType{%
%     \pgfkeysvalueof{/TitrePedagogique/type}%
%   }%

%   \ifx\TitreTmpType\empty
%     % Aucun type : rien
%   \else

%     \par
%     \vspace{0.15ex}%

%     {\centering
%       \fontsize{\TitreMetaTailleType}
%                {\TitreMetaInterligneType}
%       \selectfont
%       \itshape
%       \textcolor{#1}{\TitreTmpType}%
%       \par
%     }%

%   \fi
% }


% % ============================================================
% % TITRE PÉDAGOGIQUE GÉNÉRIQUE
% % ============================================================

% % #1 = couleur
% % #2 = taille du titre
% % #3 = largeur maximale du titre
% % #4 = épaisseur des barres
% % #5 = distance entre barre et titre
% % #6 = espace vertical après le titre
% % #7 = titre

% % \newcommand{\TitreNiveauUnOld}[1]{%
% %   \noindent\begin{tikzpicture}
% %     \node[inner sep=0pt, align=center, font=\fontsize{15}{17}\selectfont\bfseries,
% %       color=colorStitre] (title) at (0.5\linewidth,0) {
% %       \begin{varwidth}{0.85\linewidth}\centering #1\end{varwidth}
% %     };
% %     \draw[colorStitre, line width=0.3cm] (0,0) -- ([xshift=-0.5cm]title.west);
% %     \draw[colorStitre, line width=0.3cm] ([xshift=0.5cm]title.east) -- (\linewidth,0);
% %   \end{tikzpicture}
% %   \par\vspace{0.35cm}
% % }

% \newcommand{\TitrePedagogiqueAvecMeta}[7]{%

%   % ----------------------------------------------------------
%   % PARAMÈTRES VISUELS
%   % ----------------------------------------------------------
%   %
%   % Ces paramètres peuvent être changés ici globalement.
%   %
%   % Ils sont volontairement indépendants de la taille du titre.

%   \renewcommand{\TitreMetaTailleEtoile}{0.78em}
%   \renewcommand{\TitreMetaEchelleEtoile}{0.78}

%   \renewcommand{\TitreMetaTailleDuree}{8pt}
%   \renewcommand{\TitreMetaInterligneDuree}{9.5pt}

%   \renewcommand{\TitreMetaTailleType}{7.5pt}
%   \renewcommand{\TitreMetaInterligneType}{9pt}


%   % ----------------------------------------------------------
%   % DÉBUT DU TITRE
%   % ----------------------------------------------------------

%   \noindent
%   \begin{tikzpicture}[
%     baseline=(title.base)
%   ]

%     % --------------------------------------------------------
%     % TITRE PRINCIPAL
%     % --------------------------------------------------------
%     %
%     % Les métadonnées apparaissent AVANT le titre sur la même
%     % ligne.
%     %
%     % Exemple :
%     %
%     %   → ★★☆☆☆ • 25 min   Exercice 1. ...
%     %
%     % Le titre et toutes les métadonnées utilisent #1.

%     \node[
%       inner sep=0pt,
%       outer sep=0pt,
%       align=center,
%       text width=#3,
%       font=#2,
%       text=#1,
%       % draw=ellipse,
%       % opacity=0.5,
%     ] (title) at (0.5\linewidth,0) {%
%       %
%       % Métadonnées
%       %
%       {\fontsize{8pt}{9.5pt}\selectfont
%         \TitreMetaInline{#1}%
%       }%
%       %
%       % Petit espace avant le titre
%       %
%       \hspace{0.15em}%
%       %
%       % Titre
%       %
%       #7%
%     };


%     % --------------------------------------------------------
%     % TYPE D'EXERCICE
%     % --------------------------------------------------------
%     %
%     % Il est placé sous le titre et centré.
%     %
%     % Exemple :
%     %
%     %                     Application directe du cours

%     \node[
%       below=0.25ex of title.south,
%       inner sep=0pt,
%       outer sep=0pt,
%       text width=#3,
%       align=center,
%       % draw=ellipse,
%       % opacity=0.5,
%     ] (type) {%
%       \TitreMetaType{#1}%
%     };


%     % --------------------------------------------------------
%     % BARRE GAUCHE
%     % --------------------------------------------------------
%     %
%     % Épaisseur conservée de la version originale.

%     \draw[
%       color=#1,
%       line width=#4,
%     ]
%       (0,0)
%       --
%       ([xshift=-#5]title.west);


%     % --------------------------------------------------------
%     % BARRE DROITE
%     % --------------------------------------------------------

%     \draw[
%       color=#1,
%       line width=#4,
%     ]
%       ([xshift=#5]title.east)
%       --
%       (\linewidth,0);

%   \end{tikzpicture}

%   \par
%   \vspace{#6}
% }


% % ============================================================
% % NIVEAU 1
% % ============================================================

% \NewDocumentCommand{\TitreNiveauUn}{O{}m}{%

%   % ----------------------------------------------------------
%   % Valeurs par défaut
%   % ----------------------------------------------------------

%   \pgfkeys{
%     /TitrePedagogique/indispensable=false,
%     /TitrePedagogique/difficulte=,
%     /TitrePedagogique/duree=,
%     /TitrePedagogique/type=
%   }%

%   % ----------------------------------------------------------
%   % Valeurs fournies par l'utilisateur
%   % ----------------------------------------------------------

%   \pgfkeys{
%     /TitrePedagogique/.cd,
%     #1
%   }%

%   % ----------------------------------------------------------
%   % Rendu
%   %
%   % Épaisseur originale : 0.3 cm
%   % ----------------------------------------------------------

%   \TitrePedagogiqueAvecMeta
%     {colorStitre}
%     {\fontsize{15}{17}\selectfont\bfseries}
%     {0.85\linewidth}
%     {0.3cm}
%     {0.5cm}
%     {0.35cm}
%     {#2}%
% }


% % ============================================================
% % NIVEAU 2
% % ============================================================

% \NewDocumentCommand{\TitreNiveauDeux}{O{}m}{%

%   \pgfkeys{
%     /TitrePedagogique/indispensable=false,
%     /TitrePedagogique/difficulte=,
%     /TitrePedagogique/duree=,
%     /TitrePedagogique/type=
%   }%

%   \pgfkeys{
%     /TitrePedagogique/.cd,
%     #1
%   }%

%   % Épaisseur originale : 0.15 cm

%   \TitrePedagogiqueAvecMeta
%     {colorSStitre}
%     {\fontsize{12}{14}\selectfont\bfseries}
%     {0.85\linewidth}
%     {0.15cm}
%     {0.3cm}
%     {0.30cm}
%     {#2}%
% }


% % ============================================================
% % NIVEAU 3
% % ============================================================

% \NewDocumentCommand{\TitreNiveauTrois}{O{}m}{%

%   \pgfkeys{
%     /TitrePedagogique/indispensable=false,
%     /TitrePedagogique/difficulte=,
%     /TitrePedagogique/duree=,
%     /TitrePedagogique/type=
%   }%

%   \pgfkeys{
%     /TitrePedagogique/.cd,
%     #1
%   }%

%   % Épaisseur originale : 0.10 cm

%   \TitrePedagogiqueAvecMeta
%     {colorSSStitre}
%     {\fontsize{10}{12}\selectfont\bfseries}
%     {0.85\linewidth}
%     {0.10cm}
%     {0.2cm}
%     {0.25cm}
%     {#2}%
% }


% % ============================================================
% % NIVEAU 4
% % ============================================================

% \NewDocumentCommand{\TitreNiveauQuatre}{O{}m}{%

%   \pgfkeys{
%     /TitrePedagogique/indispensable=false,
%     /TitrePedagogique/difficulte=,
%     /TitrePedagogique/duree=,
%     /TitrePedagogique/type=
%   }%

%   \pgfkeys{
%     /TitrePedagogique/.cd,
%     #1
%   }%

%   % Épaisseur originale : 0.06 cm

%   \TitrePedagogiqueAvecMeta
%     {colorSSStitre}
%     {\fontsize{10}{12}\selectfont\itshape}
%     {0.85\linewidth}
%     {0.06cm}
%     {0.15cm}
%     {0.20cm}
%     {#2}%
% }


% % ============================================================
% % NIVEAU 5
% % ============================================================

% \NewDocumentCommand{\TitreNiveauCinq}{O{}m}{%

%   \pgfkeys{
%     /TitrePedagogique/indispensable=false,
%     /TitrePedagogique/difficulte=,
%     /TitrePedagogique/duree=,
%     /TitrePedagogique/type=
%   }%

%   \pgfkeys{
%     /TitrePedagogique/.cd,
%     #1
%   }%

%   % Épaisseur originale : 0.04 cm

%   \TitrePedagogiqueAvecMeta
%     {grayBrume}
%     {\fontsize{9}{11}\selectfont\itshape}
%     {0.85\linewidth}
%     {0.04cm}
%     {0.10cm}
%     {0.15cm}
%     {#2}%
% }


% % ------------------------------------------------------------------
% % Barres de titre à 5 niveaux (section -> sous-paragraphe), reprises du
% % Dossier scientifique pour les 3 premiers niveaux, complétées dans le
% % même esprit pour les 2 niveaux les plus fins (non fournis à l'origine).
% %
% % Volontairement PAS branchées sur \titleformat{\section}/... : nos
% % documents utilisent une numérotation "maison" (Exercice N, Section N)
% % gérée par les données, pas par les compteurs LaTeX -- brancher sur
% % \titleformat créerait un double numérotage. \TitreNiveauUn{...}
% % s'utilise directement à la place d'un \subsection*{...}.
% % ------------------------------------------------------------------

% % ============================================================
% % ÉTOILES DE DIFFICULTÉ
% % ============================================================

% \newcommand{\EtoilesDifficulte}[1]{%
%   \begin{tikzpicture}[
%     baseline=-0.35ex,
%     etoile/.style={
%       inner sep=0pt,
%       outer sep=0pt,
%       font=\fontsize{7.5}{7.5}\selectfont
%     }
%   ]
%     \foreach \i in {1,...,5}{%
%       \ifnum\i>#1
%         \node[
%           etoile,
%           text=grayBrume
%         ] at ({(\i-1)*0.85em},0) {$\star$};
%       \else
%         \node[
%           etoile,
%           text=colorOchre
%         ] at ({(\i-1)*0.85em},0) {$\star$};
%       \fi
%     }%
%   \end{tikzpicture}%
% }

% % ============================================================
% % MÉTADONNÉES D'EXERCICE
% % ============================================================

% \newcommand{\MetaExercice}[4]{%
%   % #1 = indispensable (true/false)
%   % #2 = difficulté
%   % #3 = durée
%   % #4 = type d'exercice
%   %
%   \begin{tikzpicture}[
%     baseline,
%     meta/.style={
%       inner sep=0pt,
%       outer sep=0pt,
%       font=\fontsize{7.5}{9}\selectfont
%     }
%   ]

%     % --------------------------------------------------------
%     % Ligne supérieure : indispensable + difficulté + durée
%     % --------------------------------------------------------

%     \ifstrequal{#1}{true}{%
%       % Flèche indispensable
%       \node[
%         meta,
%         text=colorStitre
%       ] (indispensable) {$\boldsymbol{\rightarrow}$};

%       \node[
%         meta,
%         right=0.25em of indispensable,
%         text=colorStitre,
%         font=\fontsize{7.5}{9}\selectfont\bfseries
%       ] (important) {indispensable};
%     }{%
%       \node[meta] (important) {};
%     }

%     % Difficulté
%     \node[
%       meta,
%       right=0.8em of important
%     ] (difficulte) {%
%       \EtoilesDifficulte{#2}%
%     };

%     % Durée
%     \node[
%       meta,
%       right=0.8em of difficulte,
%       text=grayPierre
%     ] (duree) {$\bullet$~#3};

%   \end{tikzpicture}

%   % ----------------------------------------------------------
%   % Type d'exercice
%   % ----------------------------------------------------------

%   \if\relax\detokenize{#4}\relax
%   \else
%     \par
%     {\color{grayBrume}%
%       \fontsize{7.5}{9}\selectfont
%       \textit{#4}%
%     }%
%   \fi
% }

% % ============================================================
% % AFFICHAGE DES MÉTADONNÉES
% % ============================================================

% \newcommand{\AfficherMetaTitre}[4]{%
%   % #1 = indispensable
%   % #2 = difficulté
%   % #3 = durée
%   % #4 = type

%   \pgfkeys{
%     /TitrePedagogique,
%     indispensable=#1,
%     difficulte=#2,
%     duree={#3},
%     type={#4}
%   }

%   % ----------------------------------------------------------
%   % Ligne des métadonnées
%   % ----------------------------------------------------------

%   \ifstrequal{\pgfkeysvalueof{/TitrePedagogique/difficulte}}{}{}{%
%     \par
%     \begin{center}
%       \begin{tikzpicture}[
%         baseline=-0.25ex
%       ]

%         \node[
%           inner sep=0pt,
%           outer sep=0pt
%         ] (stars) {%
%           \EtoilesDifficulte{
%             \pgfkeysvalueof{/TitrePedagogique/difficulte}
%           }%
%         };

%         \ifstrequal{\pgfkeysvalueof{/TitrePedagogique/duree}}{}{}{%
%           \node[
%             right=0.8em of stars,
%             inner sep=0pt,
%             outer sep=0pt,
%             font=\fontsize{7.5}{9}\selectfont,
%             text=grayPierre
%           ] (duree) {%
%             $\bullet$~\pgfkeysvalueof{/TitrePedagogique/duree}%
%           };
%         }%

%       \end{tikzpicture}
%     \end{center}
%   }%

%   % ----------------------------------------------------------
%   % Type d'exercice
%   % ----------------------------------------------------------

%   \ifstrequal{\pgfkeysvalueof{/TitrePedagogique/type}}{}{}{%
%     \par
%     \begin{center}
%       {\color{grayBrume}%
%         \fontsize{7.5}{9}\selectfont
%         \textit{\pgfkeysvalueof{/TitrePedagogique/type}}%
%       }%
%     \end{center}
%   }%
% }

% % ============================================================
% % MÉTADONNÉES DES TITRES PÉDAGOGIQUES
% % ============================================================

% \pgfkeys{
%   /TitrePedagogique/.is family,
%   /TitrePedagogique,
%   indispensable/.initial=false,
%   difficulte/.initial=,
%   duree/.initial=,
%   type/.initial=
% }


% % ============================================================
% % TITRE PÉDAGOGIQUE GÉNÉRIQUE
% % ============================================================

% \newcommand{\TitrePedagogique}[7][]{%
%   % #1 = options / métadonnées
%   % #2 = taille police
%   % #3 = interligne
%   % #4 = couleur
%   % #5 = épaisseur barre
%   % #6 = marge horizontale de la barre
%   % #7 = titre

%   \pgfkeys{
%     /TitrePedagogique,
%     indispensable=false,
%     difficulte=,
%     duree=,
%     type=
%   }

%   \pgfkeys{/TitrePedagogique,#1}

%   \noindent
%   \begin{tikzpicture}

%     % --------------------------------------------------------
%     % Titre
%     % --------------------------------------------------------

%     \node[
%       inner sep=0pt,
%       align=center,
%       font=\fontsize{#2}{#3}\selectfont\bfseries,
%       color=#4
%     ] (title) at (0.5\linewidth,0) {%
%       \begin{varwidth}{0.85\linewidth}
%         \centering
%         #7
%       \end{varwidth}
%     };

%     % --------------------------------------------------------
%     % Barre gauche
%     % --------------------------------------------------------

%     \draw[
%       color=#4,
%       line width=#5
%     ]
%       (0,0)
%       --
%       ([xshift=-#6]title.west);

%     % --------------------------------------------------------
%     % Barre droite
%     % --------------------------------------------------------

%     \draw[
%       color=#4,
%       line width=#5
%     ]
%       ([xshift=#6]title.east)
%       --
%       (\linewidth,0);

%     % --------------------------------------------------------
%     % Flèche "indispensable"
%     % --------------------------------------------------------
%     %
%     % Elle est directement dessinée dans la barre gauche.
%     %

%     \ifstrequal{\pgfkeysvalueof{/TitrePedagogique/indispensable}}{true}{%
%       \draw[
%         color=#4,
%         line width=0.8pt,
%         -{Latex[length=2.2mm,width=1.5mm]}
%       ]
%         ([xshift=0.08cm]title.west)
%         --
%         ([xshift=-0.38cm]title.west);
%     }{}%

%   \end{tikzpicture}

%   % ----------------------------------------------------------
%   % Métadonnées sous le titre
%   % ----------------------------------------------------------

%   \ifstrequal{\pgfkeysvalueof{/TitrePedagogique/difficulte}}{}{}{%
%     \AfficherMetaTitre
%       {\pgfkeysvalueof{/TitrePedagogique/indispensable}}
%       {\pgfkeysvalueof{/TitrePedagogique/difficulte}}
%       {\pgfkeysvalueof{/TitrePedagogique/duree}}
%       {\pgfkeysvalueof{/TitrePedagogique/type}}
%   }%

%   \par\vspace{0.35cm}
% }

% \newcommand{\TitreNiveauUn}[2][]{%
%   \TitrePedagogique[
%     #1
%   ]{15}{17}{colorStitre}{0.3cm}{0.5cm}{#2}%
% }

% \newcommand{\TitreNiveauDeux}[2][]{%
%   \TitrePedagogique[
%     #1
%   ]{12}{14}{colorSStitre}{0.15cm}{0.3cm}{#2}%
% }

% \newcommand{\TitreNiveauTrois}[2][]{%
%   \TitrePedagogique[
%     #1
%   ]{10}{12}{colorSSStitre}{0.10cm}{0.2cm}{#2}%
% }

% \newcommand{\TitreNiveauQuatre}[2][]{%

%   \pgfkeys{
%     /TitrePedagogique,
%     indispensable=false,
%     difficulte=,
%     duree=,
%     type=
%   }

%   \pgfkeys{/TitrePedagogique,#1}

%   \noindent
%   \begin{tikzpicture}

%     % Petite barre
%     \draw[
%       color=colorSSStitre,
%       line width=0.06cm
%     ]
%       (0,0) -- (0.4cm,0);

%     % Titre
%     \node[
%       anchor=west,
%       inner sep=0pt,
%       outer sep=0pt,
%       text=colorSSStitre,
%       font=\itshape\fontsize{10}{12}\selectfont
%     ] at (0.55cm,0) {#2};

%     % Flèche indispensable
%     \ifstrequal{\pgfkeysvalueof{/TitrePedagogique/indispensable}}{true}{%
%       \draw[
%         color=colorSSStitre,
%         line width=0.7pt,
%         -{Latex[length=2mm,width=1.4mm]}
%       ]
%         (-0.05cm,0)
%         --
%         (0.32cm,0);
%     }{}%

%   \end{tikzpicture}

%   \AfficherMetaTitre
%     {\pgfkeysvalueof{/TitrePedagogique/indispensable}}
%     {\pgfkeysvalueof{/TitrePedagogique/difficulte}}
%     {\pgfkeysvalueof{/TitrePedagogique/duree}}
%     {\pgfkeysvalueof{/TitrePedagogique/type}}

%   \par\vspace{0.2cm}
% }

% \newcommand{\TitreNiveauCinq}[2][]{%

%   \pgfkeys{
%     /TitrePedagogique,
%     indispensable=false,
%     difficulte=,
%     duree=,
%     type=
%   }

%   \pgfkeys{/TitrePedagogique,#1}

%   \noindent
%   \begin{tikzpicture}

%     % Petite barre
%     \draw[
%       color=grayBrume,
%       line width=0.04cm
%     ]
%       (0,0) -- (0.25cm,0);

%     % Titre
%     \node[
%       anchor=west,
%       inner sep=0pt,
%       outer sep=0pt,
%       text=grayBrume,
%       font=\itshape\fontsize{9}{11}\selectfont
%     ] at (0.4cm,0) {#2};

%     % Flèche indispensable
%     \ifstrequal{\pgfkeysvalueof{/TitrePedagogique/indispensable}}{true}{%
%       \draw[
%         color=grayBrume,
%         line width=0.6pt,
%         -{Latex[length=1.8mm,width=1.2mm]}
%       ]
%         (-0.05cm,0)
%         --
%         (0.18cm,0);
%     }{}%

%   \end{tikzpicture}

%   \AfficherMetaTitre
%     {\pgfkeysvalueof{/TitrePedagogique/indispensable}}
%     {\pgfkeysvalueof{/TitrePedagogique/difficulte}}
%     {\pgfkeysvalueof{/TitrePedagogique/duree}}
%     {\pgfkeysvalueof{/TitrePedagogique/type}}

%   \par\vspace{0.15cm}
% }






% \newcommand{\TitreNiveauUn}[2][]{%

%   \noindent
%   \begin{tikzpicture}

%     % --------------------------------------------------------
%     % Titre
%     % --------------------------------------------------------

%     \node[
%       inner sep=0pt,
%       align=center,
%       font=\fontsize{15}{17}\selectfont\bfseries,
%       color=colorStitre
%     ] (title) at (0.5\linewidth,0) {%
%       \begin{varwidth}{0.85\linewidth}
%         \centering
%         #2
%       \end{varwidth}
%     };

%     % --------------------------------------------------------
%     % Barres latérales
%     % --------------------------------------------------------

%     \draw[
%       color=colorStitre,
%       line width=0.3cm
%     ]
%       (0,0)
%       --
%       ([xshift=-0.5cm]title.west);

%     \draw[
%       color=colorStitre,
%       line width=0.3cm
%     ]
%       ([xshift=0.5cm]title.east)
%       --
%       (\linewidth,0);

%   \end{tikzpicture}

%   \par\vspace{0.35cm}

% }





% \newcommand{\EncadreStyle}[4]{%
%   \begin{tcolorbox}[
%       enhanced,
%       notitle,
%       colback=white,
%       colbacklower=white,
%       frame hidden,
%       boxrule=0pt,
%       sharp corners,
%       borderline west={4pt}{0pt}{#3},
%       left=7mm,
%       fontupper=\sffamily,
%       overlay={%
%           % Bande verticale
%           \fill[#3]
%           (frame.south west) rectangle
%           ([xshift=4pt]frame.north west);

%           \node[
%             rotate=90,
%             anchor=south,
%             font=\sffamily\bfseries,
%             text=#3
%           ] at ([xshift=2pt]frame.west) {#2};

%           % Texte vertical (titre)
%           \node[
%             rotate=90,
%             anchor=center,
%             font=\sffamily\bfseries\small,
%             text=white
%           ] at ([xshift=2pt]frame.west |- frame.center)
%           {#1};

%         }
%     ]
%     % contenu
%     %\textcolor{#3}{\sffamily\textbf{#1}}\\
%     #4
%   \end{tcolorbox}
% }

% ------------------------------------------------------------------

% ============================================================
% Style général des encadrés
% ============================================================

% \NewDocumentCommand{\EncadreStyle}{mmmm}{%
%   \begin{tcolorbox}[
%       enhanced,
%       notitle,
%       colback=white,
%       colbacklower=white,
%       frame hidden,
%       boxrule=0pt,
%       sharp corners,
%       borderline west={4pt}{0pt}{#3},
%       left=7mm,
%       fontupper=\sffamily,
%       overlay={%
%           % Bande verticale
%           \fill[#3]
%           (frame.south west) rectangle
%           ([xshift=4pt]frame.north west);

%           % Étiquette
%           \node[
%             rotate=90,
%             anchor=south,
%             font=\sffamily\bfseries,
%             text=#3
%           ] at ([xshift=2pt]frame.west) {#2};

%           % Titre vertical
%           \node[
%             rotate=90,
%             anchor=center,
%             font=\sffamily\bfseries\small,
%             text=white
%           ] at ([xshift=2pt]frame.west |- frame.center)
%           {#1};
%       }
%   ]
%   #4
%   \end{tcolorbox}%
% }

% \NewDocumentCommand{\EncadreStyle}{mmmm}{%
%   \def\EncadreTitre{#1}%
%   \def\EncadreEtiquette{#2}%
%   \def\EncadreCouleur{#3}%
%   \def\EncadreContenu{#4}%

%   \begin{tcolorbox}[
%     enhanced,
%     notitle,
%     colback=white,
%     colbacklower=white,
%     frame hidden,
%     boxrule=0pt,
%     sharp corners,
%     borderline west={4pt}{0pt}{\EncadreCouleur},
%     left=7mm,
%     fontupper=\sffamily,
%     overlay={%
%       % Bande verticale
%       \fill[\EncadreCouleur]
%         (frame.south west)
%         rectangle
%         ([xshift=4pt]frame.north west);

%       % Étiquette
%       \node[
%         rotate=90,
%         anchor=south,
%         font=\sffamily\bfseries,
%         text=\EncadreCouleur
%       ] at ([xshift=2pt]frame.west)
%       {\EncadreEtiquette};

%       % Titre vertical
%       \node[
%         rotate=90,
%         anchor=center,
%         font=\sffamily\bfseries\small,
%         text=black
%       ] at ([xshift=2pt]frame.west |- frame.center)
%       {\EncadreTitre};
%     }
%   ]

%   \EncadreContenu

%   \end{tcolorbox}%
% }


% ============================================================
% Environnements
% ============================================================

% \NewDocumentEnvironment{coupdepouce}{o +b}
%   {%
%     \IfNoValueTF{#1}
%       {%
%         \EncadreStyle
%           {Coup de pouce}
%           {}
%           {colorOchre}
%           {#2}%
%       }
%       {%
%         \EncadreStyle
%           {Coup de pouce}
%           {#1}
%           {colorOchre}
%           {#2}%
%       }%
%   }
%   {}


% \NewDocumentEnvironment{correction}{o +b}
%   {%
%     \IfNoValueTF{#1}
%       {%
%         \EncadreStyle
%           {Correction}
%           {}
%           {greenSapin}
%           {#2}%
%       }
%       {%
%         \EncadreStyle
%           {Correction}
%           {#1}
%           {greenSapin}
%           {#2}%
%       }%
%   }
%   {}


% \NewDocumentEnvironment{solutiondetaillee}{o +b}
%   {%
%     \IfNoValueTF{#1}
%       {%
%         \EncadreStyle
%           {Solution détaillée}
%           {}
%           {grayPierre}
%           {#2}%
%       }
%       {%
%         \EncadreStyle
%           {Solution détaillée}
%           {#1}
%           {grayPierre}
%           {#2}%
%       }%
%   }
%   {}


% \NewDocumentEnvironment{remarque}{o +b}
%   {%
%     \IfNoValueTF{#1}
%       {%
%         \EncadreStyle
%           {Remarque}
%           {}
%           {headerblue}
%           {#2}%
%       }
%       {%
%         \EncadreStyle
%           {Remarque}
%           {#1}
%           {headerblue}
%           {#2}%
%       }%
%   }
%   {}


% % ============================================================
% % Aller plus loin
% % #1 = niveau
% % ============================================================

% \NewDocumentEnvironment{allerplusloin}{o +b}
%   {%
%     \IfNoValueTF{#1}
%       {%
%         \EncadreStyle
%           {Aller plus loin}
%           {}
%           {Prune}
%           {#2}%
%       }
%       {%
%         \EncadreStyle
%           {Aller plus loin (#1)}
%           {(#1)}
%           {Prune}
%           {#2}%
%       }%
%   }
%   {}



% ============================================================
% Environnement générique
% ============================================================

% \NewDocumentEnvironment{encadrestyle}{mmm}
% {%
%   \begin{tcolorbox}[
%     %float, % <-- Rend la tcolorbox elle-même flottante
%     enhanced,
%     notitle,
%     colback=#3!5,
%     colbacklower=white,
%     frame hidden,
%     boxrule=0pt,
%     sharp corners,
%     borderline west={4pt}{0pt}{#3},
%     left=7mm,
%     fontupper=\sffamily,
%     overlay={%
%       % Bande verticale
%       \fill[#3]
%         (frame.south west)
%         rectangle
%         ([xshift=4pt]frame.north west);

%       % Étiquette
%       \node[
%         rotate=90,
%         anchor=south,
%         font=\sffamily\bfseries\fontsize{5pt}{1pt}\selectfont\bfseries,
%         text=#3,
%         inner sep=0pt
%       ] at ([xshift=-3pt]frame.west)
%       {#2};

%       % Titre
%       \node[
%         rotate=90,
%         anchor=center,
%         font=\sffamily\bfseries\small\fontsize{5pt}{10pt}\selectfont\bfseries,
%         text=white,
%         inner sep=0pt
%       ] at ([xshift=2pt]frame.west |- frame.center)
%       {#1};
%     }
%   ]
% }
% {%
%   \end{tcolorbox}%
% }


% % ============================================================
% % Environnements pédagogiques
% % ============================================================


% \NewDocumentEnvironment{prerequis}{o}
% {%
%   \IfNoValueTF{#1}
%     {\begin{encadrestyle}{Prérequis}{}{colorOchre}}
%     {\begin{encadrestyle}{Prérequis}{#1}{colorOchre}}
% }
% {%
%   \end{encadrestyle}%
% }

% \NewDocumentEnvironment{objectifspedagogiquesetudiants}{o}
% {%
%   \IfNoValueTF{#1}
%     {\begin{encadrestyle}{Objectifs pédagogiques (étudiants)}{}{colorObjectifsPedagogiquesEtudiants}}
%     {\begin{encadrestyle}{Objectifs pédagogiques (étudiants)}{#1}{colorObjectifsPedagogiquesEtudiants}}
% }
% {%
%   \end{encadrestyle}%
% }

% \NewDocumentEnvironment{objectifspedagogiquesenseignants}{o}
% {%
%   \IfNoValueTF{#1}
%     {\begin{encadrestyle}{Objectifs pédagogiques (enseignants)}{}{colorObjectifsPedagogiquesEnseignants}}
%     {\begin{encadrestyle}{Objectifs pédagogiques (enseignants)}{#1}{colorObjectifsPedagogiquesEnseignants}}
% }
% {%
%   \end{encadrestyle}%
% }

% \NewDocumentEnvironment{pointcours}{o}
% {%
%   \IfNoValueTF{#1}
%     {\begin{encadrestyle}{Point de cours}{}{colorPointCours}}
%     {\begin{encadrestyle}{Point de cours}{#1}{colorPointCours}}
% }
% {%
%   \end{encadrestyle}%
% }

% \NewDocumentEnvironment{coupdepouce}{o}
% {%
%   \IfNoValueTF{#1}
%     {\begin{encadrestyle}{Coup de pouce}{}{colorOchre}}
%     {\begin{encadrestyle}{Coup de pouce}{#1}{colorOchre}}
% }
% {%
%   \end{encadrestyle}%
% }


% \NewDocumentEnvironment{correction}{o}
% {%
%   \IfNoValueTF{#1}
%     {\begin{encadrestyle}{Correction}{}{greenSapin}}
%     {\begin{encadrestyle}{Correction}{#1}{greenSapin}}
% }
% {%
%   \end{encadrestyle}%
% }


% \NewDocumentEnvironment{solutiondetaillee}{o}
% {%
%   \IfNoValueTF{#1}
%     {\begin{encadrestyle}{Solution détaillée}{}{grayPierre}}
%     {\begin{encadrestyle}{Solution détaillée}{#1}{grayPierre}}
% }
% {%
%   \end{encadrestyle}%
% }


% \NewDocumentEnvironment{remarque}{o}
% {%
%   \IfNoValueTF{#1}
%     {\begin{encadrestyle}{Remarque}{}{headerblue}}
%     {\begin{encadrestyle}{Remarque}{#1}{headerblue}}
% }
% {%
%   \end{encadrestyle}%
% }

% \NewDocumentEnvironment{remarquespedagogiquesecueilsclassiques}{o}
% {%
%   \IfNoValueTF{#1}
%     {\begin{encadrestyle}{Remarques pédagogiques (écueils classiques)}{}{colorRemarquesPedagogiquesEcueilsClassiques}}
%     {\begin{encadrestyle}{Remarques pédagogiques (écueils classiques)}{#1}{colorRemarquesPedagogiquesEcueilsClassiques}}
% }
% {%
%   \end{encadrestyle}%
% }

% \NewDocumentEnvironment{remarquespedagogiquesconseilsanimation}{o}
% {%
%   \IfNoValueTF{#1}
%     {\begin{encadrestyle}{Remarques pédagogiques (conseils d'animation)}{}{colorRemarquesPedagogiquesConseilsAnimation}}
%     {\begin{encadrestyle}{Remarques pédagogiques (conseils d'animation)}{#1}{colorRemarquesPedagogiquesConseilsAnimation}}
% }
% {%
%   \end{encadrestyle}%
% }

% \NewDocumentEnvironment{bilan}{o}
% {%
%   \IfNoValueTF{#1}
%     {\begin{encadrestyle}{Bilan}{}{colorBilan}}
%     {\begin{encadrestyle}{Bilan}{#1}{colorBilan}}
% }
% {%
%   \end{encadrestyle}%
% }


% \NewDocumentEnvironment{allerplusloin}{o}
% {%
%   \IfNoValueTF{#1}
%     {\begin{encadrestyle}{Aller plus loin}{}{Prune}}
%     {\begin{encadrestyle}{Aller plus loin}{#1}{Prune}}
% }
% {%
%   \end{encadrestyle}%
% }



% ------------------------------------------------------------------

% ------------------------------------------------------------------
% Boîtes pédagogiques du projet : bande colorée à gauche + petite
% étiquette colorée AU-DESSUS (pas de texte pivoté à l'intérieur de la
% boîte) -- contrairement à \EncadreUn/\EncadreDeux, qui pivotent le
% titre DANS la bande, ces boîtes sont souvent courtes (1-2 lignes) et
% le texte pivoté s'est avéré illisible/coupé en pratique (testé et
% corrigé pendant le développement -- voir le rendu avant/après dans le
% message de livraison). Le rendu ci-dessous reste robuste quelle que
% soit la hauteur du contenu.
% ------------------------------------------------------------------
% \newcommand{\etiquettebox}[2]{%
%   {\color{#1}\rule{3pt}{10pt}}\,\textbf{\textsf{\small\color{#1}#2}}\\[3pt]%
% }

% \newenvironment{coupdepouce}
%   {\etiquettebox{colorOchre}{Coup de pouce}%
%    \begin{tcolorbox}[enhanced, notitle, colback=colorOchre!6, frame hidden,
%      boxrule=0pt, sharp corners, borderline west={3pt}{0pt}{colorOchre},
%      left=8pt, right=8pt, top=5pt, bottom=5pt]}
%   {\end{tcolorbox}}

% \newenvironment{correction}
%   {\etiquettebox{greenSapin}{Correction}%
%    \begin{tcolorbox}[enhanced, notitle, colback=greenSapin!5, frame hidden,
%      boxrule=0pt, sharp corners, borderline west={3pt}{0pt}{greenSapin},
%      left=8pt, right=8pt, top=5pt, bottom=5pt]}
%   {\end{tcolorbox}}

% \newenvironment{solutiondetaillee}
%   {\etiquettebox{grayPierre}{Solution détaillée}%
%    \begin{tcolorbox}[enhanced, notitle, colback=grayPierre!5, frame hidden,
%      boxrule=0pt, sharp corners, borderline west={3pt}{0pt}{grayPierre},
%      left=8pt, right=8pt, top=5pt, bottom=5pt]}
%   {\end{tcolorbox}}

% \newenvironment{remarque}
%   {\etiquettebox{headerblue}{Remarque}%
%    \begin{tcolorbox}[enhanced, notitle, colback=headerblue!5, frame hidden,
%      boxrule=0pt, sharp corners, borderline west={3pt}{0pt}{headerblue},
%      left=8pt, right=8pt, top=5pt, bottom=5pt]}
%   {\end{tcolorbox}}

% % allerplusloin garde son paramètre (niveau), affiché entre parenthèses
% % après le label principal.
% \newenvironment{allerplusloin}[1]
%   {\etiquettebox{Prune}{Aller plus loin (#1)}%
%    \begin{tcolorbox}[enhanced, notitle, colback=Prune!4, frame hidden,
%      boxrule=0pt, sharp corners, borderline west={3pt}{0pt}{Prune},
%      left=8pt, right=8pt, top=5pt, bottom=5pt]}
%   {\end{tcolorbox}}
